\subsection{Instance query complexity and input influence}\label{sec:instance}

The norm threshold concerns the worst network in a class. A useful
bound for a particular network would also recognize cancellation and
the direction in which its inputs affect the output.  The following
examples show why both are necessary.

\begin{proposition}[Equal weight magnitudes, different sampling costs]
\label{prop:magnitude-obstruction}
Let \(D\ge2\), let \(m=2^k\le n\) with \(m\ge2\), and let
\(a=A2^{-r}>0\), where \(A,r\) are nonnegative integers.  Suppose
\(b\ge r+k+1\).  There are two width-\(n\), depth-\(D\) networks
with zero biases, row norms at most \(a\), and identical matrices
\((|w_{\ell,ij}|)\) at every layer, whose outputs are respectively
\[
  0\qquad\text{and}\qquad
  F_{a,D}\left(\frac1m\sum_{i=1}^m x_i\right).
\]
The first output can be sampled with no input probes and one fair bit.
Every exact sampler for the second satisfies the lower
bound~\eqref{eq:lowerexact} and, if \(a>1\), the lower
bound~\eqref{eq:lowersuper}, with this \(m\).
\end{proposition}

\paragraph{Proof idea.}
Duplicate one signal along two paths and subtract them at the
output. Negating one path changes cancellation into addition while preserving
every absolute weight. The amplified mean-chain lower bound separates the
two sampling costs.

\begin{proof}
Only neurons 1 and 2 can be nonzero before the output layer.
In both networks the first row of the first layer has weights
\(a/m\) on the first \(m\) input coordinates.  In the first network,
the second row is identical.  In the second network, every weight
in that row is negated.  Every other first-layer entry is zero.
For \(2\le\ell\le D-1\), set
\(w_{\ell,11}=w_{\ell,22}=a\), and set all other weights to zero.
At layer \(D\), set
\[
  w_{D,11}=a/2,\qquad w_{D,12}=-a/2,
\]
and again set every other entry to zero.  When \(D=2\), the
intermediate range of layers is empty.

Write \(u=m^{-1}\sum_{i=1}^m x_i\).  In the first network the two
active neurons at layer \(D-1\) both equal \(F_{a,D-1}(u)\), so
the last preactivation is zero.  In the second network their values
are \(F_{a,D-1}(u)\) and \(-F_{a,D-1}(u)\), because tanh is odd.
The last preactivation is therefore \(aF_{a,D-1}(u)\).
This proves the two output identities.

The construction changes signs only.  Its nonzero row norms are
\(a\), and its parameters belong to the grid \(2^{-b}\mathbb Z\):
the only divisions of \(a\) are by \(m=2^k\) and by 2.
The mean-chain lower bound in Theorem~\ref{thm:lower}, which holds
for every power of two \(m\le n\) by
Remark~\ref{rem:lower-general-m}, applies unchanged to the second
network, regardless of the preprocessing.
\end{proof}

For \(D=\lceil\log_2n\rceil\) and fixed dyadic \(a>1\), this
is a polynomial separation.  Thus row norms, absolute path products,
and even the complete sequence of entrywise absolute weight matrices
cannot determine instance sampling cost within polylogarithmic
factors.  The output of the first network is a fair coin; a uniform
algorithm can detect this displayed cancellation during preprocessing.

\begin{proposition}[A zero product profile can conceal a hard network]
\label{prop:profile-obstruction}
Let \(m\le n\) be even, and define
\[
 f(x)=F_{a,D}\left(
 \frac1m\sum_{i=1}^{m/2}x_i-
 \frac1m\sum_{i=m/2+1}^{m}x_i\right).
\]
Whenever \(a\) and \(a/m\) belong to the parameter grid, this is
the output of a legal width-\(n\), depth-\(D\) network with row
norms at most \(a\) and zero biases.  Its optimal worst-input
expected query cost equals that of the unsigned mean chain.
Nevertheless, under independent input signs of common mean \(t\),
\[
 H_f(t)=\E_t f(X)=0\qquad(-1\le t\le1).
\]
Consequently all derivatives of this one-parameter profile vanish.
\end{proposition}

\paragraph{Proof idea.}
Flipping half the input signs preserves the query problem but
changes the direction of its hard average. Exchanging the two halves leaves
the common-mean input distribution invariant and negates the output, forcing
that one-dimensional profile to vanish.

\begin{proof}
Use weights \(a/m\) on the first half of the first row and
\(-a/m\) on its second half.  Subsequent distinguished rows
have their sole nonzero weight equal to \(a\).  All other weights
and all biases are zero.  This realizes the displayed function.

Replacing \(x_i\) by \(-x_i\) in the second half maps this
function to the unsigned mean chain.  A query to either input can
be answered by one query to its image and a deterministic sign flip.
The map is an involution, so it gives reductions in both directions
without changing the number of probes.

To compute the product profile, exchange the first and second input
halves.  The common-mean product law is invariant under this
permutation.  The argument of \(F_{a,D}\) is negated, and the
function \(F_{a,D}\) is odd.  Hence \(\E_t f(X)=-\E_t f(X)\).
The profile is zero.  It is a polynomial in \(t\), so every
derivative is also zero.
\end{proof}

This example rules out every instance statistic that uses only the
single profile \(t\mapsto\E_t f(X)\), including all of its
derivatives.  It leaves the multivariate product extension available.

\subsubsection{A multivariate information certificate}

For \(\boldsymbol t\in(-1,1)^m\), let
\[
 \mathcal H(\boldsymbol t)=
 \sum_{x\in\{-1,1\}^m} f(x)
 \prod_{i=1}^m\frac{1+t_ix_i}{2}.
\]
Let \(q_i(\boldsymbol t)\) be the probability that an arbitrary
exact sampler probes coordinate \(i\) under this product law.
Repeated queries can be answered from a cache.  For a direction
\(v\in\mathbb R^m\), put
\[
 I_v(\boldsymbol t)=\sum_{i=1}^m
 \frac{v_i^2}{1-t_i^2}\,q_i(\boldsymbol t).
\]

\begin{lemma}[Directional transcript inequalities]
\label{lem:multivariate-transcript}
Every almost surely terminating exact sign sampler satisfies
\begin{align}
 \bigl(v\cdot\nabla\mathcal H(\boldsymbol t)\bigr)^2
 &\le (1-\mathcal H(\boldsymbol t)^2)I_v(\boldsymbol t),
 \label{eq:multivariate-information}\\
 \left|v^{\mathsf T}\nabla^2\mathcal H(\boldsymbol t)v\right|
 &\le2I_v(\boldsymbol t).
 \label{eq:multivariate-curvature}
\end{align}
The assertions allow arbitrary weight-only preprocessing and require
no moment bound on arithmetic time.
\end{lemma}

\paragraph{Proof idea.}
Perturb all input biases in a chosen direction and differentiate
the transcript likelihood. Fresh-coordinate scores are orthogonal martingale
increments. Their variance controls the first derivative, and the second
likelihood derivative controls curvature.

\begin{proof}
Follow the finite decision-tree reduction used in
Lemma~\ref{lem:transcript}.  On a neighborhood of zero, perturb the
input mean vector to \(\boldsymbol t+uv\).  At \(u=0\), the
first and negative second log-likelihood derivatives of a leaf are
\[
 Z=\sum_{i\text{ queried}}\frac{v_iX_i}{1+t_iX_i},
 \qquad
 A=\sum_{i\text{ queried}}\frac{v_i^2}{(1+t_iX_i)^2}.
\]
A new coordinate retains its product law conditional on the past.
Its score has conditional mean zero and conditional second moment
\(v_i^2/(1-t_i^2)\).  The same expression is the conditional
mean of its contribution to \(A\).  There are at most \(m\)
distinct probes, and therefore
\[
 \E Z=0,\qquad \E Z^2=\E A=I_v(\boldsymbol t).
\]
Finite likelihood differentiation gives
\[
 v\cdot\nabla\mathcal H=\E[YZ],\qquad
 v^{\mathsf T}\nabla^2\mathcal H v=\E[Y(Z^2-A)].
\]
Since \(\E Z=0\), the first expectation equals
\(\E[(Y-\mathcal H)Z]\).  Cauchy--Schwarz and
\(\operatorname{Var}(Y)=1-\mathcal H^2\) prove
\eqref{eq:multivariate-information}.  The inequalities
\(|Y|=1\) and \(|Z^2-A|\le Z^2+A\) prove
\eqref{eq:multivariate-curvature}.
\end{proof}

For a nonconstant \(f\), the product extension lies in \((-1,1)\)
at every interior point.  Taking directions with
\(v_i^2\le1-t_i^2\) gives the instance lower bound
\begin{equation}\label{eq:instance-certificate}
 \sup_x\E[Q\mid x]\ge
 \sup_{\boldsymbol t\in(-1,1)^m}
 \max\left\{
 \frac{\left(\sum_i\sqrt{1-t_i^2}
                  |\partial_i\mathcal H(\boldsymbol t)|\right)^2}
      {1-\mathcal H(\boldsymbol t)^2},
 \frac12\sup_{v_i^2\le1-t_i^2}
 \left|v^{\mathsf T}\nabla^2\mathcal H(\boldsymbol t)v\right|
 \right\}.
\end{equation}
The first term follows by choosing each \(v_i\) with the sign of
\(\partial_i\mathcal H\).  The second term follows directly
from \eqref{eq:multivariate-curvature}.  If \(f\) is constant,
both terms are defined to be zero.

This certificate recognizes the signed mean chain: take
\(t_i=t\) in its positive half and \(t_i=-t\) in its negative
half, then differentiate in the corresponding signed direction.
Computing the displayed extension by its definition entails
\(2^m\) network evaluations.
